\documentclass[10pt,a4paper]{amsart}
\usepackage{amsmath,amssymb,amsthm,graphicx}
\usepackage{verbatim}
\newenvironment{reference}{\comment}{\endcomment}
%\newenvironment{reference}{}{}
\newenvironment{slogan}{\comment}{\endcomment}
\newenvironment{history}{\comment}{\endcomment}

% For commutative diagrams we use Xy-pic
\usepackage[all]{xy}

% We use 2cell for 2-commutative diagrams.
\xyoption{2cell}
% Theorem environments.
%
\theoremstyle{plain}
\newtheorem{theorem}[subsection]{Theorem}
\newtheorem{proposition}[subsection]{Proposition}
\newtheorem{lemma}[subsection]{Lemma}

\theoremstyle{definition}
\newtheorem{definition}[subsection]{Definition}
\newtheorem{example}[subsection]{Example}
\newtheorem{exercise}[subsection]{Exercise}
\newtheorem{situation}[subsection]{Situation}

\theoremstyle{remark}
\newtheorem{remark}[subsection]{Remark}
\newtheorem{remarks}[subsection]{Remarks}

\numberwithin{equation}{subsection}

% Macros
%
\def\lim{\mathop{\mathrm{lim}}\nolimits}
\def\colim{\mathop{\mathrm{colim}}\nolimits}
\def\Spec{\mathop{\mathrm{Spec}}}
\def\Hom{\mathop{\mathrm{Hom}}\nolimits}
\def\Ext{\mathop{\mathrm{Ext}}\nolimits}
\def\SheafHom{\mathop{\mathcal{H}\!\mathit{om}}\nolimits}
\def\SheafExt{\mathop{\mathcal{E}\!\mathit{xt}}\nolimits}
\def\Sch{\mathit{Sch}}
\def\Mor{\mathop{\mathrm{Mor}}\nolimits}
\def\Ob{\mathop{\mathrm{Ob}}\nolimits}
\def\Sh{\mathop{\mathit{Sh}}\nolimits}
\def\NL{\mathop{N\!L}\nolimits}
\def\CH{\mathop{\mathrm{CH}}\nolimits}
\def\proetale{{pro\text{-}\acute{e}tale}}
\def\etale{{\acute{e}tale}}
\def\QCoh{\mathit{QCoh}}
\def\Ker{\mathop{\mathrm{Ker}}}
\def\Im{\mathop{\mathrm{Im}}}
\def\Coker{\mathop{\mathrm{Coker}}}
\def\Coim{\mathop{\mathrm{Coim}}}

% Boxtimes
%
\DeclareMathSymbol{\boxtimes}{\mathbin}{AMSa}{"02}

%
% Macros for moduli stacks/spaces
%
\def\QCohstack{\mathcal{QC}\!\mathit{oh}}
\def\Cohstack{\mathcal{C}\!\mathit{oh}}
\def\Spacesstack{\mathcal{S}\!\mathit{paces}}
\def\Quotfunctor{\mathrm{Quot}}
\def\Hilbfunctor{\mathrm{Hilb}}
\def\Curvesstack{\mathcal{C}\!\mathit{urves}}
\def\Polarizedstack{\mathcal{P}\!\mathit{olarized}}
\def\Complexesstack{\mathcal{C}\!\mathit{omplexes}}
% \Pic is the operator that assigns to X its picard group, usage \Pic(X)
% \Picardstack_{X/B} denotes the Picard stack of X over B
% \Picardfunctor_{X/B} denotes the Picard functor of X over B
\def\Pic{\mathop{\mathrm{Pic}}\nolimits}
\def\Picardstack{\mathcal{P}\!\mathit{ic}}
\def\Picardfunctor{\mathrm{Pic}}
\def\Deformationcategory{\mathcal{D}\!\mathit{ef}}
\usepackage{cleveref}
\crefname{originallemma}{original lemma}{original lemmas}
\Crefname{originallemma}{Original lemma}{Original lemmas}
\makeatletter
\newlabel{descent-lemma-syntomic-smooth-etale-permanence}{{02KM}{source}{}{sourcepointer-1}{}}
\newlabel{descent-lemma-syntomic-smooth-etale-permanence@cref}{{[originallemma][0][]02KM}{[1][0][]source}{}{}{}}
\makeatother
\crefname{originallemma}{original lemma}{original lemmas}
\Crefname{originallemma}{Original lemma}{Original lemmas}
\makeatletter
\newlabel{morphisms-lemma-composition-etale}{{02GN}{source}{}{sourcepointer-2}{}}
\newlabel{morphisms-lemma-composition-etale@cref}{{[originallemma][0][]02GN}{[1][0][]source}{}{}{}}
\makeatother
\crefname{originallemma}{original lemma}{original lemmas}
\Crefname{originallemma}{Original lemma}{Original lemmas}
\makeatletter
\newlabel{morphisms-lemma-etale-permanence}{{02GW}{source}{}{sourcepointer-3}{}}
\newlabel{morphisms-lemma-etale-permanence@cref}{{[originallemma][0][]02GW}{[1][0][]source}{}{}{}}
\makeatother
\crefname{originaldefinition}{original definition}{original definitions}
\Crefname{originaldefinition}{Original definition}{Original definitions}
\makeatletter
\newlabel{spaces-definition-relative-representable-property}{{025V}{source}{}{sourcepointer-4}{}}
\newlabel{spaces-definition-relative-representable-property@cref}{{[originaldefinition][0][]025V}{[1][0][]source}{}{}{}}
\makeatother
\crefname{originallemma}{original lemma}{original lemmas}
\Crefname{originallemma}{Original lemma}{Original lemmas}
\makeatletter
\newlabel{spaces-lemma-composition-representable-transformations-property}{{02WK}{source}{}{sourcepointer-5}{}}
\newlabel{spaces-lemma-composition-representable-transformations-property@cref}{{[originallemma][0][]02WK}{[1][0][]source}{}{}{}}
\makeatother
\crefname{originallemma}{original lemma}{original lemmas}
\Crefname{originallemma}{Original lemma}{Original lemmas}
\makeatletter
\newlabel{spaces-lemma-morphism-schemes-gives-representable-transformation-property}{{02WJ}{source}{}{sourcepointer-6}{}}
\newlabel{spaces-lemma-morphism-schemes-gives-representable-transformation-property@cref}{{[originallemma][0][]02WJ}{[1][0][]source}{}{}{}}
\makeatother
\title{Composition, cancellation and descent of etale morphisms into an algebraic space}
\author{The Stacks Project Authors}
\date{stacks-project:spaces-properties.tex:local-snapshot-20261012; GNU FDL 1.2 or later}
\begin{document}
\maketitle
\noindent\textit{Editorial note.} All delivered mathematical text is the author's own, copied from the pinned Stacks Project chapter file spaces-properties.tex; no statement or proof step was rewritten, reordered or removed.\par\medskip
\noindent\textit{Editorial note.} The counted proof is the author's complete argument: the first part through the definition and stability under composition, the second part by base change along a surjective etale chart with descent of etaleness, and the third part by base change along an arbitrary scheme over the space and a further descent step.\par\medskip
\noindent\textit{Editorial note.} The label lemma-etale-over-space names the counted statement, and the proof environment immediately following it in the source is delivered verbatim as its proof.\par\medskip
\noindent\textit{Editorial note.} The counted statement is carried by the author's own lemma environment in the chapter source.\par\medskip
\noindent\textit{Editorial note.} Pre-open gate: python operations/agent-authoring/tools/precard\_gate.py --source spaces-properties.tex --label lemma-etale-over-space --start 2236 --end 2309 --local-id 28049 returned RESULT=OK (L28049 free, occupancy FREE, public interval clear). Candidate taken from sources/stacks-spaces-properties/candidates-halfA.json.\par\medskip
\noindent\textit{Editorial note.} The statement and proof are copied from the pinned chapter file; the cited inputs are kept as external pointers with their canonical Stacks tags.\par\medskip
\noindent\textit{Allowed standard prerequisites (editorial).}\begin{itemize}
\item Definition of an etale morphism from a scheme into an algebraic space. Spaces, the definition of relative representable properties, together with the fact that any morphism from a scheme into an algebraic space is representable. Supplies the interpretation of etaleness into an algebraic space as etaleness after base change to any scheme.
\item Permanence of etaleness under composition and base change. Morphisms, the lemmas that etale morphisms are stable under composition and under base change. Supplies the etaleness of the composites and of the base changes appearing in the three parts.
\item Descent of etaleness along a surjective etale morphism. Descent, the lemma allowing the etale property of a morphism of schemes to be checked after a surjective etale base change. Supplies the conclusion in parts (2) and (3) that the morphism in question is etale.
\end{itemize}\medskip
\ifdefined\hypertarget\hypertarget{sourcepointer-1}{}\fi
\noindent\textit{Source reference.} Original lemma 02KM: sources/stacks-chow/tags.txt, source lines 3377--3377. Prerequisite pointer invoked at proof line(s) 2291; it names a standard fact of a sibling chapter that the counted argument uses, and every step of the counted proof built on it is retained verbatim in the delivered body.\par\medskip
\ifdefined\hypertarget\hypertarget{sourcepointer-2}{}\fi
\noindent\textit{Source reference.} Original lemma 02GN: sources/stacks-chow/tags.txt, source lines 3234--3234. Prerequisite pointer invoked at proof line(s) 2268; it names a standard fact of a sibling chapter that the counted argument uses, and every step of the counted proof built on it is retained verbatim in the delivered body.\par\medskip
\ifdefined\hypertarget\hypertarget{sourcepointer-3}{}\fi
\noindent\textit{Source reference.} Original lemma 02GW: sources/stacks-chow/tags.txt, source lines 3243--3243. Prerequisite pointer invoked at proof line(s) 2281; it names a standard fact of a sibling chapter that the counted argument uses, and every step of the counted proof built on it is retained verbatim in the delivered body.\par\medskip
\ifdefined\hypertarget\hypertarget{sourcepointer-4}{}\fi
\noindent\textit{Source reference.} Original definition 025V: sources/stacks-chow/tags.txt, source lines 2828--2828. Prerequisite pointer invoked at proof line(s) 2260; it names a standard fact of a sibling chapter that the counted argument uses, and every step of the counted proof built on it is retained verbatim in the delivered body.\par\medskip
\ifdefined\hypertarget\hypertarget{sourcepointer-5}{}\fi
\noindent\textit{Source reference.} Original lemma 02WK: sources/stacks-chow/tags.txt, source lines 3807--3807. Prerequisite pointer invoked at proof line(s) 2271; it names a standard fact of a sibling chapter that the counted argument uses, and every step of the counted proof built on it is retained verbatim in the delivered body.\par\medskip
\ifdefined\hypertarget\hypertarget{sourcepointer-6}{}\fi
\noindent\textit{Source reference.} Original lemma 02WJ: sources/stacks-chow/tags.txt, source lines 3806--3806. Prerequisite pointer invoked at proof line(s) 2272; it names a standard fact of a sibling chapter that the counted argument uses, and every step of the counted proof built on it is retained verbatim in the delivered body.\par\medskip
\begin{lemma}
\label{lemma-etale-over-space}
Let $S$ be a scheme.
Let $X$ be an algebraic space over $S$.
Let $U$, $U'$ be schemes over $S$.
\begin{enumerate}
\item If $U \to U'$ is an \'etale morphism of schemes, and
if $U' \to X$ is an \'etale morphism from $U'$ to $X$, then the
composition $U \to X$ is an \'etale morphism from $U$ to $X$.
\item If $\varphi : U \to X$ and $\varphi' : U' \to X$ are
\'etale morphisms towards $X$, and if $\chi : U \to U'$ is a
morphism of schemes such that $\varphi = \varphi' \circ \chi$,
then $\chi$ is an \'etale morphism of schemes.
\item If $\chi : U \to U'$ is a surjective \'etale morphism
of schemes and $\varphi' : U' \to X$ is a morphism such that
$\varphi = \varphi' \circ \chi$ is \'etale, then $\varphi'$
is \'etale.
\end{enumerate}
\end{lemma}

\begin{proof}
Recall that our definition of an \'etale morphism from a scheme into an
algebraic space comes from
Spaces, Definition
\ref{spaces-definition-relative-representable-property}
via the fact that any morphism from a scheme into an algebraic space
is representable.

\medskip\noindent
Part (1) of the lemma follows from this, the fact that
\'etale morphisms are preserved under composition
(Morphisms, Lemma
\ref{morphisms-lemma-composition-etale})
and
Spaces, Lemmas
\ref{spaces-lemma-composition-representable-transformations-property} and
\ref{spaces-lemma-morphism-schemes-gives-representable-transformation-property}
(which are formal).

\medskip\noindent
To prove part (2) choose a scheme $W$ over $S$ and a
surjective \'etale morphism $W \to X$. Consider the base change
$\chi_W : W \times_X U \to W \times_X U'$ of $\chi$.
As $W \times_X U$ and $W \times_X U'$ are \'etale over $W$, we conclude that
$\chi_W$ is \'etale, by
Morphisms, Lemma \ref{morphisms-lemma-etale-permanence}.
On the other hand, in the commutative diagram
$$
\xymatrix{
W \times_X U \ar[r] \ar[d] & W \times_X U' \ar[d] \\
U \ar[r] & U'
}
$$
the two vertical arrows are \'etale and surjective.
Hence by
Descent, Lemma \ref{descent-lemma-syntomic-smooth-etale-permanence}
we conclude that $U \to U'$ is \'etale.

\medskip\noindent
To prove part (3) choose a scheme $W$ over $S$ and a morphism $W \to X$.
As above we consider the diagram
$$
\xymatrix{
W \times_X U \ar[r] \ar[d] & W \times_X U' \ar[d] \ar[r] & W \ar[d] \\
U \ar[r] & U' \ar[r] & X
}
$$
Now we know that $W \times_X U \to W \times_X U'$ is surjective \'etale
(as a base change of $U \to U'$)
and that $W \times_X U \to W$ is \'etale. Thus $W \times_X U' \to W$
is \'etale by Descent, Lemma
\ref{descent-lemma-syntomic-smooth-etale-permanence}. By definition
this means that $\varphi'$ is \'etale.
\end{proof}

\end{document}
